\uuid{FqJP}
\titre{Limite en un point}
\auteur{bodin ; hanani ; mzari}
\organisation{exo7}
\licence{CC BY-NC-SA 4.0}
\datecreate{2020-01-20}
\niveau{L1}
\module{Analyse}
\chapitre{Continuité, limite et étude de fonctions réelles}
\sousChapitre{Limite de fonctions}
\difficulte{1}
\theme{Limites des fonctions réelles, limite en un point}
\exo7qcmid{lille-1.limites.003}

\contenu{
\begin{question}
\qclassification{123.03}
\qtags{motcle=limite en un point}
Soit $f(x)= \frac{1}{x+1}+ \frac{3x}{(x+1)(x^2-x+1)}$. Quelles sont les assertions vraies ?
\begin{answers}
    \bad{$\lim_{x\to -1^+} f(x)=+\infty$}

    \bad{$\lim_{x\to -1^-} f(x)=-\infty$}

    \good{$\lim_{x\to -1} f(x)=0$}

    \bad{$\lim_{x\to -1} f(x)=-2$}
\end{answers}
\begin{explanations}
En réduisant au même dénominateur, on obtient :  $f(x)= \frac{x+1}{x^2-x+1}$.
\end{explanations}
\end{question}
}
